\documentclass[10pt]{article}
\usepackage{icmlko}

\icmlmeta{IP 파운데이션 월드모델}{61}{라벨을 빼자 절반이 돌아왔다}{신경망이 진 이유가 구조도 용량도 아니었다 --- 표현을 라벨로 학습한 것이 격차의 61퍼센트다}

\begin{document}

\twocolumn[
\icmlheader

\begin{abstract}
노트 60은 공유 인코더가 현행 선형 파이프라인에 크게 진다고 보고하고
($+$0.1702 대 $+$0.3506) 반증 조건을 적었다 --- ``잠재 차원을 둘로 묶은 것이
불리했을 수 있다.'' 검정했다. \textbf{무관하다.} 잠재를 2에서 16까지 늘려도
$+$0.108에서 $+$0.119 사이로 평평하고, 층을 하나 더 쌓으면 오히려 나빠진다
($+$0.170$\to+$0.111). 용량이 문제가 아니다. 그러면 무엇인가. 도메인별 주성분
분석과 공유 인코더의 결정적 차이가 하나 남는다 --- \textbf{주성분은 라벨을 안
쓰고 인코더는 쓴다.} 그리고 노트 38 · 40이 이미 그 법칙을 찾아 두었다: 자기
라벨에 잘 맞출수록 남에게 못 옮겨진다. 표현 학습 층위에서 같은 일이 벌어진다면,
인코더를 \emph{비지도}로 학습하면 회복돼야 한다. 오토인코더로 잠재를 만들고
같은 프로크루스테스 정렬을 붙였다. \textbf{$+$0.2004에서 $+$0.3099로 올랐다.}
격차 0.180을 셋으로 쪼개면 --- 정렬을 되돌려 받은 몫 0.030(17\%), \textbf{라벨을
뺀 몫 0.110(61\%)}, 남은 표현 차이 0.041(23\%). 신경망이 진 주된 이유는
\textbf{표현을 라벨로 학습한 것}이다. 프론티어 모델이 자기지도로 사전학습하는
이유가 여섯 도메인 2,181건에서 정량으로 재현된다.
\end{abstract}
]

\section{용량이 아니었다}

노트 60의 반증 조건부터 처리했다.

\begin{table}[t]
\centering\tsize
\caption{잠재 차원과 폭을 늘릴 때. 씨앗 둘, 세 층 인코더.}
\begin{tabular}{rrrr}
\toprule
잠재 & 폭 & 평균 $\rho$ & 씨앗 SD \\
\midrule
2 & 32 & $+$0.1106 & 0.0114 \\
4 & 32 & $+$0.1127 & 0.0133 \\
8 & 64 & $+$0.1187 & 0.0115 \\
16 & 64 & $+$0.1081 & 0.0327 \\
\bottomrule
\end{tabular}
\end{table}

평평하다. 그리고 세 층으로 깊게 한 이 표가 노트 60의 두 층($+$0.1702)보다
전반적으로 나쁘다 --- 표본이 작아 깊이가 해롭다.

\textbf{용량을 늘리는 방향은 막혀 있다.} 반증 조건이 발동하지 않았으므로
노트 60의 결론은 그대로 서 있고, 원인을 다른 데서 찾아야 한다.

\section{남은 차이는 라벨이었다}

도메인별 주성분 분석과 공유 인코더의 차이를 하나씩 지웠는데 아직 안 지운 것이
있었다.

\begin{quote}\fontsize{9.5}{11.5}\selectfont
\textbf{주성분은 라벨을 쓰지 않는다.} 분산 구조만 본다. 인코더는 회귀 손실로
학습하므로 학습 도메인들의 라벨에 맞춘 표현을 만든다.
\end{quote}

노트 38과 40이 그 상황의 법칙을 이미 찾아 두었다 --- 도메인의 인자 공간이 자기
라벨에 잘 맞을수록 대상으로서 좋아지고 \emph{출처로서 나빠진다}($r=-0.85$에서
$-0.99$). 인코더는 학습 도메인들의 라벨에 특수화되므로, 그 법칙이 표현 학습
층위에서 반복된다면 비지도로 바꾸면 회복돼야 한다.

오토인코더로 잠재를 만들었다 --- 같은 입력 10차원, 같은 잠재 2차원, 재구성
손실만. 그리고 노트 60과 같은 프로크루스테스 정렬을 붙였다.

\begin{figure}[t]
\centering
\includegraphics[width=\textwidth]{figs/ladder.pdf}
\caption{왼쪽 --- 한 단계씩 되돌릴 때. 오른쪽 --- 격차의 분해.}
\end{figure}

\begin{table}[t]
\centering\tsize
\caption{격차를 단계로 되돌린다.}
\begin{tabular}{lr}
\toprule
구성 & 평균 $\rho$ \\
\midrule
지도 인코더 헤드 직접 & $+$0.1702 \\
\quad $+$ 프로크루스테스 정렬 & $+$0.2004 \\
\quad $+$ 라벨 제거(오토인코더) & \textbf{$+$0.3099} \\
도메인별 PCA(현행) & $+$0.3506 \\
\bottomrule
\end{tabular}
\end{table}

\begin{table}[t]
\centering\tsize
\caption{격차 0.180의 분해.}
\begin{tabular}{lrr}
\toprule
원인 & 몫 & 비율 \\
\midrule
\textbf{라벨로 표현을 학습한 것} & \textbf{0.110} & \textbf{61\%} \\
남은 표현 차이(PCA 대 오토인코더) & 0.041 & 23\% \\
정렬을 학습으로 대체한 것 & 0.030 & 17\% \\
\bottomrule
\end{tabular}
\end{table}

\claimx{신경망이 진 주된 이유는 구조도 용량도 정렬도 아니고 \textbf{표현을
라벨로 학습한 것}이다. 라벨을 빼면 격차의 61\%가 회복된다. 노트 38 · 40의
상충 법칙이 표현 학습 층위에서 반복된다.}{오토인코더 잠재가 여전히 PCA에
0.041 뒤진다. 그 몫이 비선형 때문인지 최적화 때문인지 안 봤다. 선형
오토인코더로 재면 갈린다.}

\section{프론티어가 왜 그렇게 하는가}

큰 모델들이 자기지도로 사전학습하고 라벨은 마지막에 얇은 헤드로만 쓴다.
그 설계의 근거가 여기서 정량으로 나온다 --- \textbf{라벨로 표현을 빚으면 그
라벨이 나온 분포에 표현이 묶인다.}

이 프로젝트에서는 그것이 특히 아프다. 도메인이 여섯인데 각자 다른 시장이고,
한 시장의 라벨에 맞춘 표현은 다른 시장에서 쓸모가 준다. 반대로 축의 분산
구조는 시장을 건너 비슷하다 --- 노트 33이 ``전이 상관 $=$ 대상 자기 상관''을
얻은 것도 결국 대상의 \emph{비지도} 구조가 전이를 정한다는 뜻이었다.

\begin{quote}\fontsize{9.5}{11.5}\selectfont
현행 파이프라인은 우연히 옳은 설계였다. 도메인별 주성분(비지도) $+$ 프로크루스테스
정렬(닫힌 해) $+$ 출처 회귀(얇은 헤드)는 \textbf{자기지도 사전학습 뒤 선형
프로빙}과 같은 구조다. 다만 사전학습이 학습이 아니라 계산이다.
\end{quote}

\section{한계}

\begin{itemize}[leftmargin=1.1em,itemsep=0.15em]
\item 오토인코더 잠재도 PCA에 0.041 뒤진다. 비선형이 이 자료에서 얻을 것이
      없다는 뜻일 수 있고, 학습이 덜 됐다는 뜻일 수도 있다.
\item 재구성 손실 하나만 썼다. 대조 학습이나 마스킹 복원 같은 다른 자기지도
      목적은 안 봤다.
\item 씨앗 셋으로 쟀다. 오토인코더 씨앗 SD가 0.0039로 작아 충분해 보이지만
      짝지은 구간은 안 붙였다.
\item ``상충 법칙이 표현 층위에서 반복된다''는 해석이지 직접 검정이 아니다.
      인코더 학습에 쓴 도메인 수를 바꿔 가며 재면 검정된다.
\item 여전히 현행 선형이 최선이다. 신경망 쪽에서 얻은 것은 설명이지 성능이 아니다.
\end{itemize}

\section{다음}

\begin{enumerate}[leftmargin=1.3em,itemsep=0.15em]
\item 선형 오토인코더로 재 남은 0.041이 비선형 탓인지 가른다.
\item 마스킹 복원 --- 축 하나를 가리고 맞히게 하는 자기지도. 결측이 많은 이
      자료에 맞는다.
\item 인코더 학습 도메인 수를 바꿔 상충 법칙 해석을 직접 검정한다.
\item 웹툰 전체 수집 완료 후 전면 재측정(1,621/3,973 진행 중).
\end{enumerate}

\vspace{0.6em}
\hrule height 0.3pt
\vspace{0.5em}
{\footnotesize
\textbf{참고문헌}\\[0.3em]
Caruana, R. (1997). Multitask learning. \emph{Machine Learning}, 28(1), 41--75.\\[0.15em]
Grinsztajn, L., Oyallon, E., Varoquaux, G. (2022). Why do tree-based models still
outperform deep learning on typical tabular data? \emph{NeurIPS D\&B}.\\[0.15em]
Chen, T. et al. (2020). A simple framework for contrastive learning of visual
representations. \emph{ICML}, 1597--1607.\\[0.15em]
Ben-David, S. et al. (2010). A theory of learning from different domains.
\emph{Machine Learning}, 79(1--2), 151--175.\\[0.15em]
Schönemann, P. H. (1966). A generalized solution of the orthogonal Procrustes
problem. \emph{Psychometrika}, 31(1), 1--10.
}

\end{document}
